%=====================================================================
%  CADENCE — IEEE conference paper (IEEEtran, conference mode)
%  Compile on Overleaf:  pdfLaTeX -> BibTeX -> pdfLaTeX x2
%  Every quantitative claim traces to a measured entry in the project's
%  results ledger (docs/results.md).  Unrun experiments are stated as
%  limitations, never as results.
%=====================================================================
\documentclass[conference]{IEEEtran}
\IEEEoverridecommandlockouts

\usepackage{cite}
\usepackage{amsmath,amssymb,amsfonts}
\usepackage{algorithmic}
\usepackage{graphicx}
\usepackage{textcomp}
\usepackage{xcolor}
\usepackage{booktabs}
\usepackage{array}
\usepackage{multirow}
\usepackage{pifont}
\usepackage{url}
\usepackage[hidelinks]{hyperref}

\graphicspath{{figures/}}

% compact check/cross marks for the comparison matrix
\newcommand{\yes}{\textcolor{black}{\ding{51}}}
\newcommand{\no}{\textcolor{gray}{\ding{55}}}

\begin{document}

\title{CADENCE: Causal Attribution and Cost-Constrained\\
Retraining-Scope Optimization for Continually\\ Deployed Machine Learning Systems}

\author{%
\IEEEauthorblockN{Gona Lalith}
\IEEEauthorblockA{Department of Computer Science Engineering\\
Lovely Professional University\\
Phagwara, India\\
gonalalith2005@gmail.com}
\and
\IEEEauthorblockN{Ankita Maji}
\IEEEauthorblockA{Department of Computer Science Engineering\\
Lovely Professional University\\
Phagwara, India\\
ankitamaji7033@gmail.com}
\and
\IEEEauthorblockN{Nihar Kanakam}
\IEEEauthorblockA{Department of Computer Science Engineering\\
Lovely Professional University\\
Phagwara, India\\
niharkanakam@gmail.com}
\and
\IEEEauthorblockN{Quaise Mallick}
\IEEEauthorblockA{Department of Computer Science Engineering\\
Lovely Professional University\\
Phagwara, India\\
quaisemallick.wha@gmail.com}
\and
\IEEEauthorblockN{Shailaja Singh}
\IEEEauthorblockA{Department of Computer Science Engineering\\
Lovely Professional University\\
Phagwara, India\\
sshailaja2004@gmail.com}
\and
\IEEEauthorblockN{Jasmeen Khatun}
\IEEEauthorblockA{Department of Computer Science Engineering\\
Lovely Professional University\\
Phagwara, India\\
jasmeenkhatun536@gmail.com}
\and
\IEEEauthorblockN{Dr.\ Ravin Kumar}
\IEEEauthorblockA{School of Computer Science and Engineering\\
Lovely Professional University\\
Punjab, India\\
ravinpal2009@gmail.com}
}

\maketitle

\begin{abstract}
Deployed machine-learning models degrade as their input and label
distributions drift. The two dominant industrial responses---retrain
on a fixed schedule, or retrain fully whenever a drift alarm fires---are
both blunt: they ignore \emph{why} the model degraded and pay for a
full retrain even when a small, targeted repair would suffice. We present
CADENCE, a closed-loop framework that (i) builds a \emph{Causal Drift
Attribution Graph} (CDAG) spanning both external input features and
internal activation clusters, (ii) ranks the responsible nodes with a
graph neural network and a counterfactual surrogate that predicts the
benefit of repairing each node, and (iii) selects a minimal retraining
scope with a cost-constrained reinforcement-learning policy formalized as
a constrained Markov decision process (CMDP). We evaluate CADENCE under a
\emph{pre-registered} protocol that fixes success criteria before running,
and we report negative results as first-class findings. On a synthetic
benchmark with ground-truth root causes, the CDAG+GNN attribution beats a
Population-Stability-Index baseline on the contested (non-saturated)
scenarios by a wide margin (mean reciprocal rank $0.10{\to}0.41$;
AUROC $0.67{\to}0.91$; paired Wilcoxon $p<0.001$, $n{=}10$ seeds). On a
genuinely multi-task benchmark (Split-MNIST), an EWC-regularized partial
retrain reduces catastrophic forgetting from $39.5\%$ to near zero
against both a naive-full and a fair full-with-replay baseline
($p<0.001$, $n{=}10$). On real electricity-market drift (Elec2), the
policy cuts $69.1\%$ of retraining compute for a $1.7$-point F1 cost
($95\%$ CI $[-0.025,-0.009]$)---a Pareto-frontier operating point rather
than a strict dominance. Accounting for CADENCE's own overhead (dominated by
causal-graph construction) we find---and report---that this net saving holds
only when the avoided retrain is expensive (measured break-even
$\approx1.3$\,s). We formalize the constrained objective, provide a
falsifiable diagnostic that the learned policy is healthy (non-collapsed,
exploring, constraint-satisfying), and delineate exactly which claims are
established and which await a pre-registered confirmatory run.
\end{abstract}

\begin{IEEEkeywords}
concept drift, causal attribution, graph neural networks, constrained
reinforcement learning, continual learning, MLOps, model monitoring.
\end{IEEEkeywords}

%=====================================================================
\section{Introduction}
%=====================================================================
A model that is accurate at deployment does not stay accurate. As the
serving distribution drifts away from the training distribution, quality
silently erodes~\cite{gama2014survey,lu2018learning,quinonero2009dataset}.
The operational cost of this is well documented: technical debt
accumulates around models that no one can safely change, and the majority
of the effort in a deployed ML system is spent \emph{after} the model is
first trained~\cite{sculley2015debt,paleyes2022challenges}.

Two responses dominate practice. \emph{Periodic} retraining refits the
model on a fixed calendar, paying for compute whether or not anything
changed. \emph{Reactive} retraining fires a full refit whenever a
statistical drift alarm (e.g.\ Population Stability Index, PSI, or a
Kolmogorov--Smirnov test) crosses a threshold. Both are blunt in the same
way: they treat the model as a black box and answer only \emph{when} to
retrain, never \emph{why it broke} or \emph{how much of it to repair}. A
covariate shift in a single feature and a wholesale change in the labeling
concept demand very different responses, yet both trigger the same
expensive full retrain.

We argue that the right unit of decision is \emph{minimum-cost repair
under a performance constraint}: given a degradation, identify the
responsible part of the input--model system, estimate the benefit of
repairing it, and execute the smallest intervention that restores service
within a service-level agreement (SLA). This reframes model maintenance as
a causal-attribution problem feeding a constrained sequential-decision
problem.

\textbf{Contributions.} This paper makes four contributions, each
supported by measured evidence and each with its scope stated honestly:
\begin{enumerate}
\item \textbf{A hybrid causal drift-attribution graph (CDAG)} whose nodes
span both external input features \emph{and} internal activation clusters,
with a learned GNN readout. On the contested scenarios where correlational
PSI cannot resolve the true root cause, the learned scorer wins on rank
and ranking-AUROC metrics with $p<0.001$ (Sec.~\ref{sec:h1}).
\item \textbf{A counterfactual surrogate} that predicts the F1 benefit of
repairing each candidate node, trained on sandbox interventions and
calibrated on held-out data ($R^2=0.89$), together with an honest analysis
of where it fails (out-of-distribution, unrecoverable drift; Sec.~\ref{sec:surrogate}).
\item \textbf{A cost-constrained Retraining-Scope Optimizer (RSO)}
formalized as a CMDP and solved with an augmented-Lagrangian PPO. We
provide a \emph{pre-registered diagnostic} establishing that the learned
policy is non-degenerate---non-collapsed, exploring, and
constraint-satisfying---which is a prerequisite that many RL-for-systems
papers omit (Sec.~\ref{sec:rso},~\ref{sec:diagnostic}).
\item \textbf{A pre-registered, reproducibility-first evaluation}
spanning synthetic and real drift and three model families, in which
negative results (single-task forgetting, unrecoverable concept flips,
masked drift under weak encoders) are reported as findings rather than
hidden (Sec.~\ref{sec:results},~\ref{sec:limits}).
\end{enumerate}

We are deliberate about what is and is not established. The attribution
(H1) and forgetting (H3) claims are supported at paper scale with
$p<0.001$. The cost-saving (H2) claim is a \emph{practical} Pareto result
on one real dataset at $n{=}3$; its strict-dominance form is the subject
of a pre-registered confirmatory run whose protocol we publish
(Sec.~\ref{sec:limits}). We regard the discipline of separating the two as
a feature, not a hedge.

%=====================================================================
\section{Related Work and Positioning}\label{sec:related}
%=====================================================================
\textbf{Drift detection} is mature: distributional tests such as PSI and
KS, windowed detectors (ADWIN~\cite{bifet2007adwin}), error-rate monitors
(DDM~\cite{gama2004ddm}), and learned shift detectors~\cite{rabanser2019failing}
answer \emph{whether} the distribution moved. They do not localize a cause
inside the feature--model system, and they do not decide a repair.

\textbf{Causal discovery} provides the machinery to move from correlation
to structure: constraint-based PC~\cite{spirtes2000causation}, continuous
structure learning (NOTEARS~\cite{zheng2018notears}), and time-series
causal discovery~\cite{runge2019pcmci} (surveyed in~\cite{glymour2019review}),
grounded in the interventional
semantics of Pearl~\cite{pearl2009causality} and
Peters~et~al.~\cite{peters2017elements}. These methods target graph
\emph{recovery}; they are not tied to a downstream repair decision, and
their guarantees weaken on the short windows available at alert time
(a limitation we measure and discuss in Sec.~\ref{sec:limits}).

\textbf{Feature attribution} (SHAP~\cite{lundberg2017shap}, Integrated
Gradients~\cite{sundararajan2017ig}, permutation
importance~\cite{breiman2001random}) explains \emph{predictions} on a
fixed model, not \emph{degradation} of that model under drift, and it is
purely correlational with respect to the drift cause.

\textbf{Continual learning} addresses forgetting when a model is updated:
EWC~\cite{kirkpatrick2017ewc}, synaptic
intelligence~\cite{zenke2017si}, replay/GEM~\cite{lopezpaz2017gem}, and
exemplar methods~\cite{rebuffi2017icarl}, surveyed
in~\cite{delange2021continual}. These are \emph{mechanisms} for a repair,
not policies for choosing \emph{whether and how much} to repair.

\textbf{Adaptive/triggered retraining} and MLOps platforms
(TFX~\cite{baylor2017tfx}, production monitoring~\cite{klaise2020monitoring},
data validation~\cite{polyzotis2019data}, production-readiness
rubrics~\cite{breck2017mltest}) operationalize retraining but
typically on schedules or thresholds rather than a learned,
cost-constrained policy.

\textbf{Positioning.} CADENCE is, to our knowledge, the first framework to
place all of drift detection, causal attribution over a graph that
includes \emph{internal} model state, counterfactual benefit prediction,
a \emph{learned cost-constrained} retraining-scope policy, and
forgetting-aware execution into one closed loop. Table~\ref{tab:matrix}
makes the comparison concrete; each row cites a representative method so
the claim is falsifiable rather than asserted.

\begin{table*}[t]
\centering
\caption{Capability comparison against representative prior methods. A
mark denotes the capability is a designed, first-class part of the method,
not that it is impossible to bolt on. ``Internal nodes'' means attribution
over internal model state (e.g.\ activations), not only inputs.}
\label{tab:matrix}
\renewcommand{\arraystretch}{1.25}
\setlength{\tabcolsep}{4pt}
\begin{tabular}{@{}lccccccc@{}}
\toprule
\textbf{Method (representative)} & \textbf{Drift} & \textbf{Causal} & \textbf{Internal} & \textbf{Counterfact.} & \textbf{Retrain} & \textbf{Cost} & \textbf{Continual}\\
 & \textbf{detect} & \textbf{attrib.} & \textbf{nodes} & \textbf{interv.} & \textbf{scope} & \textbf{constraint} & \textbf{learning}\\
\midrule
PSI / KS / ADWIN~\cite{bifet2007adwin}                     & \yes & \no  & \no  & \no  & \no  & \no  & \no  \\
Learned shift detection~\cite{rabanser2019failing}         & \yes & \no  & \no  & \no  & \no  & \no  & \no  \\
Feature attribution (SHAP/IG)~\cite{lundberg2017shap,sundararajan2017ig} & \no & \no & \no & \no & \no & \no & \no \\
Causal discovery (PC/NOTEARS)~\cite{spirtes2000causation,zheng2018notears} & \no & \yes & \no & varies & \no & \no & \no \\
Continual learning (EWC/GEM)~\cite{kirkpatrick2017ewc,lopezpaz2017gem}     & \no & \no  & \no  & \no  & \yes & \no  & \yes \\
Triggered / periodic retrain~\cite{baylor2017tfx}          & \yes & \no  & \no  & \no  & varies & \no  & varies \\
\textbf{CADENCE (this work)}                               & \yes & \yes & \yes & \yes & \yes & \yes & \yes \\
\bottomrule
\end{tabular}
\end{table*}

%=====================================================================
\section{Problem Formulation}\label{sec:formulation}
%=====================================================================
\subsection{What ``causal attribution'' means here}
We are explicit to avoid over-claiming graph \emph{recovery}. Let a
degradation be observed at monitoring time $t$. Let
$\mathcal{I}=\{I_1,\dots,I_m\}$ be candidate intervention targets (input
features and internal activation clusters). For target $I_i$, define its
\emph{repair responsibility} as the interventional improvement in
performance obtained by repairing $I_i$ and holding others fixed:
\begin{equation}
R_i \;=\; \mathbb{E}\!\left[\,\mathrm{F1}\big(Y \mid do(I_i{=}\text{repaired})\big)\right] - \mathrm{F1}^{\text{cur}} .
\label{eq:responsibility}
\end{equation}
CADENCE does \emph{not} claim to recover the true causal DAG. It claims
to (B) \emph{rank} targets by $R_i$ and (C) \emph{predict} $R_i$ with a
learned surrogate $\hat R_i = g_\theta(G, X, i)$, where $G$ is the CDAG and
$X$ the current window. The scientific question is therefore how well
$\hat R_i$ orders and predicts the true $R_i$, which we can measure without
ground-truth graph recovery (Sec.~\ref{sec:h1},~\ref{sec:surrogate}). The
CDAG is a useful \emph{inductive bias} for $g_\theta$, not a claim of
identifiability on short windows.

\subsection{Retraining as a constrained MDP}
The RSO chooses, each decision step, how much of the model to repair. We
model this as a CMDP $(\mathcal{S},\mathcal{A},P,r,c,\gamma)$:
\begin{itemize}
\item \textbf{State} $s_t\in\mathbb{R}^{15}$: drift statistics (per-feature
PSI summary), current and recent F1, the SLA margin $\mathrm{SLA}-\mathrm{F1}_t$,
the attribution \emph{concentration} (how peaked $\hat R$ is across nodes),
and cost-budget context.
\item \textbf{Action} $a_t\in\{\textsc{no-op},\textsc{partial},\textsc{full}\}$:
skip, EWC-regularized partial retrain of the responsible layer(s), or a
full retrain.
\item \textbf{Transition} $P(s_{t+1}\mid s_t,a_t)$: apply the action in a
retraining sandbox, advance the drift stream one window, re-measure.
\item \textbf{Cost} $c(s_t,a_t)$: estimated GPU-hours (and carbon) of the
action.
\item \textbf{Constraint} $g(s_t,a_t)=\max(0,\ \mathrm{SLA}-\mathrm{F1}_{t+1})$:
SLA violation.
\end{itemize}
The objective is the standard CMDP form---maximize return subject to a
bound on expected violation:
\begin{equation}
\max_{\pi}\ \mathbb{E}_\pi\!\Big[\textstyle\sum_t \gamma^t\,\Delta \mathrm{F1}_t\Big]
\quad\text{s.t.}\quad \mathbb{E}_\pi\!\Big[\textstyle\sum_t \gamma^t\,c_t\Big]\le B,\ \ \mathbb{E}_\pi[g_t]\le \delta .
\label{eq:cmdp}
\end{equation}
We solve~\eqref{eq:cmdp} with an augmented-Lagrangian relaxation, giving
the per-step reward
\begin{equation}
r_t \;=\; \Delta \mathrm{F1}_t \;-\; w\,c_t \;-\; \lambda_t\,\max\!\big(0,\ \mathrm{SLA}-\mathrm{F1}_{t+1}\big),
\label{eq:reward}
\end{equation}
with the dual variable updated by projected dual ascent toward a target
violation rate $\delta$:
\begin{equation}
\lambda_{t+1} \;=\; \mathrm{clip}\big(\lambda_t + \eta_\lambda(\bar g_t - \delta),\ 0,\ \lambda_{\max}\big).
\label{eq:dual}
\end{equation}
\textbf{Why this formulation.} The cost/SLA trade-off is a constraint, not
a fixed scalarization: augmented Lagrangian~\cite{bertsekas1996constrained}
lets $\lambda$ \emph{adapt} to the realized violation rate rather than
being hand-tuned, following the constrained-RL
line~\cite{altman1999constrained,achiam2017cpo,tessler2019rcpo,chow2018lyapunov,ray2019safetygym}.
The action is sequential because each repair changes the state the next
decision faces (a full retrain now changes what drift looks like later);
a myopic classifier cannot represent that dependence, which is the standard
argument for an MDP over a one-shot predictor~\cite{sutton2018reinforcement}. We do not yet claim
PPO is \emph{necessary}---establishing that requires the ablation we
pre-register in Sec.~\ref{sec:limits}. We use PPO~\cite{schulman2017ppo}
with GAE~\cite{schulman2016gae} via
Stable-Baselines3~\cite{raffin2021sb3}.

%=====================================================================
\section{The CADENCE Architecture}\label{sec:arch}
%=====================================================================
Fig.~\ref{fig:arch} shows the closed loop. \emph{Monitor}: a
privacy-filtered collector samples input windows and internal activations
from the deployed model and runs a PSI + delayed-label F1 trigger.
\emph{Attribute}: on an alert, the CDAG builder constructs a graph over
feature nodes and per-layer activation-cluster nodes using
PC~\cite{spirtes2000causation} for the skeleton and a
NOTEARS-style~\cite{zheng2018notears} continuous relaxation for weights; a
GCN/GraphSAGE readout~\cite{kipf2017gcn,hamilton2017graphsage} (message
passing~\cite{gilmer2017mpnn}; attention variants~\cite{velickovic2018gat}
are drop-in) over the CDAG produces node embeddings consumed by the
counterfactual surrogate.
\emph{Decide}: the RSO (Sec.~\ref{sec:rso}) selects a scope.
\emph{Repair}: an executor performs an EWC-regularized
partial~\cite{kirkpatrick2017ewc} or full retrain, then a
shadow/canary validation promotes or rolls back. The loop then resumes
monitoring. Implementation uses PyTorch~\cite{paszke2019pytorch}, PyTorch
Geometric~\cite{fey2019pyg}, LightGBM~\cite{ke2017lightgbm} and
scikit-learn~\cite{pedregosa2011scikit}; cost/carbon use a hardware model
in the spirit of~\cite{lacoste2019quantifying,strubell2019energy,schwartz2020green}.

\begin{figure}[t]
\centering
\includegraphics[width=\columnwidth]{fig_architecture.pdf}
\caption{CADENCE closed loop. Monitor $\rightarrow$ attribute (CDAG +
GNN + counterfactual surrogate) $\rightarrow$ decide (cost-constrained
RSO) $\rightarrow$ repair (EWC partial / full + shadow validation)
$\rightarrow$ redeploy. The attribution graph spans external features and
internal activation clusters.}
\label{fig:arch}
\end{figure}

\subsection{Retraining-Scope Optimizer}\label{sec:rso}
The RSO is the PPO policy for~\eqref{eq:cmdp}--\eqref{eq:dual}. Two design
choices matter for reproducibility. First, we train an \emph{independent}
policy per random seed (rather than one shared policy) so cross-seed
statistics reflect genuine training variance, not a single lucky run.
Second, the dual variable $\lambda$ is updated online toward a data-informed
target violation $\delta$; Sec.~\ref{sec:diagnostic} shows this keeps
$\lambda$ off its cap, which earlier configurations did not.

%=====================================================================
\section{Experimental Setup}\label{sec:setup}
%=====================================================================
\textbf{Pre-registration.} Before running the confirmatory experiments we
fixed, in a signed document, the Stage-1 policy-health gates and the
Stage-2 verdict criteria (STRONG / PRACTICAL / INCONCLUSIVE /
NO-ADVANTAGE), following the pre-registration
rationale~\cite{nosek2018preregistration,pineau2021reproducibility}. This
prevents post-hoc metric selection and is why we can report negative
results without penalty.

\textbf{Datasets.} Synthetic drift with ground-truth root cause is
injected into a Credit-Card-Fraud MLP~\cite{dalpozzolo2015calibrating}
across five scenarios (abrupt/gradual covariate shift on two features,
an additive shift, a concept flip, and a temporal drift). Real drift:
Elec2 electricity pricing~\cite{harries1999splice} and an Airlines
delay stream. Cross-model generality: LightGBM~\cite{ke2017lightgbm} on
Give-Me-Some-Credit and a TF-IDF + logistic-regression text classifier on
Yelp reviews. Forgetting: Split-MNIST~\cite{lecun1998mnist} with disjoint
digit tasks. Optimization uses Adam~\cite{kingma2015adam}.

\textbf{Protocol.} Unless stated, results are over multiple seeds with
paired Wilcoxon signed-rank tests and $95\%$ bootstrap confidence
intervals. Every number in this paper is regenerated from a measured
artifact; no values are hand-entered.

%=====================================================================
\section{Results}\label{sec:results}
%=====================================================================

\subsection{H1: causal attribution beats correlation on contested drift}\label{sec:h1}
We compare three scorers---PSI (correlational), a structural CDAG
path-sum proxy, and the learned CDAG+GNN---on top-1/top-3 accuracy, mean
reciprocal rank (MRR), and ranking AUROC, over five scenarios $\times$ ten
seeds ($50$ paired samples). Reporting the pooled average would be
misleading: on the three ``easy'' scenarios PSI is already saturated at
$1.0$, leaving nothing to beat, which dilutes the effect. We therefore
report the \emph{easy} (PSI-saturated) and \emph{hard} (contested) subsets
separately (Fig.~\ref{fig:h1}, Table~\ref{tab:h1}).

On the \textbf{hard} subset the CDAG+GNN scorer wins decisively: MRR rises
from $0.097$ (PSI) to $0.412$ (a $4.3\times$ lift) and AUROC from $0.672$
(near chance on a 30-feature space) to $0.905$, with paired Wilcoxon
$p<0.001$ vs PSI on both metrics and $p=0.032$ vs the structural proxy.
On the \textbf{easy} subset PSI is at $1.0$ and the GNN is
statistically not-inferior. Notably the structural proxy is \emph{worse}
than PSI---evidence that propagating NOTEARS weights alone is
insufficient; the learned readout is doing the work.

\begin{figure}[t]
\centering
\includegraphics[width=\columnwidth]{fig_h1_subset.pdf}
\caption{H1 at paper scale ($n{=}10$ seeds). Solid bars: EASY
(PSI-saturated) scenarios; hatched: HARD (contested). On HARD the
CDAG+GNN beats both baselines on MRR and AUROC ($p<0.001$ vs PSI); on EASY
PSI is saturated at $1.0$.}
\label{fig:h1}
\end{figure}

\begin{table}[t]
\centering
\caption{H1 attribution, HARD subset (2 contested scenarios $\times$ 10
seeds, mean$\pm$std). Paired Wilcoxon vs the learned scorer.}
\label{tab:h1}
\setlength{\tabcolsep}{4pt}
\begin{tabular}{@{}lcccc@{}}
\toprule
\textbf{Scorer} & \textbf{top-1} & \textbf{MRR} & \textbf{AUROC} & \textbf{$p$ (MRR)}\\
\midrule
PSI (correlational)      & $0.00$          & $0.097{\pm}.01$ & $0.672{\pm}.05$ & $<0.001$\\
CDAG structural proxy    & $0.00$          & $0.150{\pm}.10$ & $0.586{\pm}.34$ & $0.032$\\
\textbf{CDAG + GNN}      & $\mathbf{0.25}$ & $\mathbf{0.412{\pm}.34}$ & $\mathbf{0.905{\pm}.06}$ & ---\\
\bottomrule
\end{tabular}
\end{table}

\subsection{Counterfactual surrogate accuracy}\label{sec:surrogate}
The surrogate $\hat R_i=g_\theta(G,X,i)$ is trained jointly with the GNN on
sandbox intervention triples $(\text{node},\text{intervention},\text{post-fix F1})$.
On the held-out validation split it is well-calibrated:
$R^2=0.888$, MAE $=0.036$ F1 points (Fig.~\ref{fig:surrogate}). Its
failure mode is instructive and honestly reported: on a genuinely
\emph{unrecoverable} concept flip---out of the surrogate's training
distribution---it over-predicts recovery (episode MAE $0.57$), defaulting
to its ``the intervention should help'' prior. This is precisely the case
the unrecoverable-drift guard is designed to catch.

\textbf{Sample-efficiency (does it amortize?).} Because each sandbox
intervention is itself a retrain, the surrogate is only worthwhile if a
\emph{few} interventions suffice. We swept the training budget: at $10$
triples the surrogate is useless ($R^2\!\approx\!0$), but it reaches
$R^2\!\approx\!0.75$ at $25$, $R^2\!\approx\!0.9$ (MAE $0.02$) by $100$, and
plateaus $\approx0.94$ thereafter---so $\sim\!50$--$100$ interventions are
enough, cheap against the wrong-retrain decisions it averts. This is the
amortization claim, measured.

\begin{figure}[t]
\centering
\includegraphics[width=0.92\columnwidth]{surrogate_calibration.png}
\caption{Counterfactual surrogate: predicted vs measured post-repair F1 on
the held-out split ($R^2=0.89$, MAE $=0.036$). The surrogate amortizes
sandbox interventions into a cheap benefit estimate.}
\label{fig:surrogate}
\end{figure}

\subsection{H2: cost-constrained repair on real drift}\label{sec:h2}
On Elec2, the serving stream degrades from a train F1 of $0.84$ to an
un-adapted stream F1 of $0.43$---a genuine $41$-point drift, not an
artifact (verified: temporal split, no leakage, class balance shifts
$-2.9$pp). Against a reactive-full baseline, CADENCE saves \textbf{69.1\%}
of retraining GPU-hours for a \textbf{$-1.66$-point} F1 gap ($95\%$
bootstrap CI $[-0.025,-0.009]$); against a periodic baseline it wins
$+0.16$ F1 at $1.54\times$ cost (Fig.~\ref{fig:pareto}). CADENCE therefore
sits on the Pareto frontier: neither baseline dominates it, and it
dominates neither strictly.

We are precise about the status of this result. Under our pre-registration
it is a \textbf{PRACTICAL} (Pareto) result, not a \textbf{STRONG}
(strict-dominance) one: the compute saving is exact and reproducible, the
F1 gap is small with a wholly-negative CI, but $n{=}3$ seeds and a
deterministic rule policy make the paired test underpowered for a
non-inferiority claim. The STRONG form requires the $n{\ge}10$, learned-policy
confirmatory run whose full protocol we publish (Sec.~\ref{sec:limits}).
On the synthetic sweep, a naive short-horizon PPO instead learns
bang-bang policies---an honest negative that directly motivated the
diagnostic in Sec.~\ref{sec:diagnostic}. Crucially, this figure is a
\emph{retraining-compute} saving; Sec.~\ref{sec:overhead} accounts for
CADENCE's own overhead and identifies exactly when it becomes a net
system saving.

\begin{figure}[t]
\centering
\includegraphics[width=0.92\columnwidth]{fig_elec2_pareto.pdf}
\caption{Real-drift Elec2 ($n{=}3$). CADENCE (rule policy) sits on the
cost/F1 Pareto frontier: $-69.1\%$ GPU-hours vs reactive-full for a
$-0.017$ F1 gap (CI $[-0.025,-0.009]$).}
\label{fig:pareto}
\end{figure}

\subsection{Total system overhead: when is the saving \emph{net}?}\label{sec:overhead}
A cost-aware framework must account for its \emph{own} compute, or the
headline saving is an illusion. We measured every component's wall-clock on
the RTX~4070 over a representative $1024$-row window
(Table~\ref{tab:overhead}) and composed the net budget using the measured
Elec2 action profile ($N{=}10$ windows; CADENCE $\{5,4,3\}$ no-op/partial/full;
reactive-full $\{0,0,10\}$):
\begin{align}
C_{\text{cad}} &= N(t_{\text{mon}}{+}t_{\text{cdag}}{+}t_{\text{gnn}}{+}t_{\text{sur}}{+}t_{\text{pol}}) + n_p t_p + n_f t_f,\nonumber\\
C_{\text{rf}} &= N\,t_{\text{mon}} + N\,t_f .
\end{align}
The overhead is dominated by CDAG construction (PC~+~NOTEARS) at
$374$\,ms/window; GNN, surrogate, monitor and policy are all $\le5$\,ms
combined. On the \emph{tiny} FraudNet a full retrain costs only $0.31$\,s,
so the composed budget is $C_{\text{cad}}{=}5.50$\,s vs
$C_{\text{rf}}{=}3.08$\,s: \textbf{CADENCE spends $78\%$ \emph{more} total
compute than reactive-full}---the $44.7\%$ retraining saving is eaten by the
attribution overhead. We report this openly. Solving $C_{\text{rf}}>C_{\text{cad}}$
for this action profile gives a break-even full-retrain cost of
$\approx\!1.3$\,s: below it, CADENCE's overhead dominates; above it---i.e.\ on
any realistically-sized production model, where a retrain takes
seconds-to-minutes---the $374$\,ms CDAG cost is negligible and the retraining
saving is essentially the net saving. \textbf{The value proposition is
therefore regime-dependent: CADENCE pays off precisely when the retrain it
avoids is expensive.} This also argues for building the CDAG only on a
trigger alert and caching discovery. We regard measuring this boundary---and
finding our own method net-negative in the cheap-retrain regime---as more
useful than a saving quoted without its overhead.

\begin{table}[t]
\centering
\caption{Measured component wall-times (RTX~4070, median of 25 reps) and the
composed net budget over the Elec2 action profile. CDAG dominates overhead;
break-even full-retrain cost $\approx1.3$\,s.}
\label{tab:overhead}
\setlength{\tabcolsep}{5pt}
\begin{tabular}{@{}lr@{}}
\toprule
\textbf{Component (per window)} & \textbf{cost}\\
\midrule
monitor (PSI)              & $1.49$\,ms\\
CDAG (PC + NOTEARS)        & $\mathbf{374.1}$\,ms\\
GNN forward               & $4.92$\,ms\\
surrogate forward         & $0.23$\,ms\\
policy                    & $0.007$\,ms\\
\midrule
partial retrain (2 ep.)   & $0.194$\,s\\
full retrain (5 ep.)      & $0.307$\,s\\
\midrule
$C_{\text{cadence}}$ (10 win.) & $5.50$\,s\\
$C_{\text{reactive-full}}$     & $3.08$\,s\\
\textbf{net vs reactive-full} & $\mathbf{-78.5\%}$\\
\bottomrule
\end{tabular}
\end{table}

\subsection{H3: targeted partial retraining reduces forgetting}\label{sec:h3}
On Split-MNIST (task A $=\{0,1\}$, task B $=\{2,3\}$, genuinely disjoint),
we compare EWC-regularized partial retraining against a naive full retrain
and a reviewer-fair full-with-replay baseline ($n{=}10$ seeds,
Fig.~\ref{fig:h3}, Table~\ref{tab:h3}). Partial retraining preserves task
A almost perfectly (forgetting $-0.03$, i.e.\ a slight gain), while the
naive full retrain loses $39.5\%$ of task-A F1; even the fair
full-with-replay baseline is matched or beaten. Both paired Wilcoxon tests
reach the $n{=}10$ theoretical floor, $p<0.001$---partial wins on
\emph{every} seed. The honest trade-off is visible in the task-B column:
EWC buys stability at a plasticity cost (task-B F1 $0.80$ vs $0.99$),
exactly as the stability--plasticity dilemma
predicts~\cite{mccloskey1989catastrophic,french1999catastrophic,delange2021continual}.

Crucially, the same hypothesis \emph{failed} on single-task
Credit-Card-Fraud: with no distinct task to preserve, a full retrain with
replay is already sufficient. Reporting both (Table~\ref{tab:h3} caption)
gives the correct, bounded claim: \emph{EWC-partial's
forgetting-protection materializes when there is a distinct task to
preserve}.

\textbf{Does \emph{causal targeting} drive this---an honest null.} A careful
reviewer asks whether the win comes from EWC or from \emph{which} layer the
attribution selects. We isolated the targeting axis: under identical EWC
(penalty, Fisher, $\theta^\star$, epochs, data), we chose the retrained layer
by random, PSI (input drift $\rightarrow$ input layer), gradient magnitude,
or CDAG activation-drift, sweeping the EWC penalty $\in\{0,10^2,10^3,5{\times}10^3\}$
($n{=}10$). The result is a null for the strong claim: \textbf{CDAG targeting
does \emph{not} reduce forgetting relative to random or gradient targeting}
(paired Wilcoxon $p\approx0.5$--$0.9$, i.e.\ tied). On this benchmark
forgetting is floored near zero for every target because the tasks are highly
separable; the targeting choice instead governs \emph{plasticity}
(input-layer retraining reaches task-B F1 $0.78$ vs $0.50$ for a deep layer).
The honest conclusion is that the forgetting benefit is owed to EWC and
partial retraining \emph{itself}, not to causal layer selection---on this
benchmark. A harder, confusable task pair and the learned (rather than
structural) scorer are the natural stronger tests.

\begin{figure}[t]
\centering
\includegraphics[width=0.92\columnwidth]{fig_h3_mnist_forget.pdf}
\caption{H3 Split-MNIST task-A forgetting ($n{=}10$). Partial-EWC (ours)
preserves task A (median $\approx-0.03$) while naive-full loses $\approx0.40$
F1. Partial $<$ naive and partial $<$ full-with-replay, both $p<0.001$.}
\label{fig:h3}
\end{figure}

\begin{table}[t]
\centering
\caption{H3 on Split-MNIST ($n{=}10$). Forgetting = task-A F1 drop after
learning task B (lower is better; negative = gain). On single-task Fraud
the effect reverses---full-with-replay suffices---which we report as a
scoping negative.}
\label{tab:h3}
\setlength{\tabcolsep}{5pt}
\begin{tabular}{@{}lccc@{}}
\toprule
\textbf{Retrain arm} & \textbf{Task-A forget} & \textbf{Task-B F1} & \textbf{$p$ vs ours}\\
\midrule
Naive full (no replay)   & $+0.395{\pm}.02$ & $0.990$ & $<0.001$\\
Full + replay (fair)     & $-0.017{\pm}.00$ & $0.979$ & $<0.001$\\
\textbf{Partial + EWC (ours)} & $\mathbf{-0.032{\pm}.00}$ & $0.802$ & ---\\
\bottomrule
\end{tabular}
\end{table}

\subsection{Cross-model generality}\label{sec:generality}
The identical decision loop runs unchanged over a neural MLP, a LightGBM
tree ensemble, and a TF-IDF/logistic-regression text model. We state the
scope precisely (Table~\ref{tab:generality}): the \emph{architecture} is
model-agnostic, but the \emph{partial-repair mechanism} is not---tree and
linear adapters expose no internal layers, so the executor's
\texttt{partial\_fit} raises \texttt{NotImplementedError} and the ladder
falls back to a full retrain by contract. On GMSC at a realistic SLA, the
tree adapter still places CADENCE on the Pareto frontier
(Fig.~\ref{fig:treepareto}): it saves $50\%$ of retrain compute vs
reactive-full for a $0.007$ F1 cost.

\begin{table}[t]
\centering
\caption{Cross-model generality. Attribution and full repair are
universal; partial repair needs internal layers.}
\label{tab:generality}
\setlength{\tabcolsep}{6pt}
\begin{tabular}{@{}lccc@{}}
\toprule
\textbf{Model family} & \textbf{Attribution} & \textbf{Partial repair} & \textbf{Full repair}\\
\midrule
MLP (neural)             & \yes & \yes & \yes\\
GBDT (LightGBM)          & \yes & \no  & \yes\\
Logistic reg.\ (TF-IDF)  & \yes & \no  & \yes\\
\bottomrule
\end{tabular}
\end{table}

\begin{figure}[t]
\centering
\includegraphics[width=0.86\columnwidth]{tree_pareto_frontier.png}
\caption{LightGBM adapter on GMSC at a realistic SLA: CADENCE saves $50\%$
of retrain compute vs reactive-full for a $0.007$ F1 cost.}
\label{fig:treepareto}
\end{figure}

\subsection{Robustness}\label{sec:robustness}
Two robustness measurements (Fig.~\ref{fig:robust}). (i)~The PSI trigger
has a $0.00$ false-alarm rate on an undrifted stream and a $1.00$
true-alert rate on a drifted one in our configuration---strong
specificity and sensitivity. (ii)~The GNN attribution degrades gracefully
under missing telemetry: randomly zeroing up to $50\%$ of CDAG edges costs
only $\sim4$ AUROC points ($0.945\to0.903$), consistent with message
passing aggregating soft signal rather than relying on any single edge.

\begin{figure}[t]
\centering
\includegraphics[width=0.9\columnwidth]{robustness_edge_dropout.png}
\caption{GNN attribution AUROC vs fraction of CDAG edges randomly dropped
(proxy for missing activation telemetry). Half the edges removed costs
$\sim4$ AUROC points.}
\label{fig:robust}
\end{figure}

\subsection{Is the learned policy healthy? A pre-registered diagnostic}\label{sec:diagnostic}
RL-for-systems results are worthless if the policy silently collapsed to a
constant action. Before any confirmatory H2 run, we therefore
\emph{pre-registered} eight Stage-1 diagnostic gates on policy health and
required them to pass. The final configuration passes seven of eight
(Table~\ref{tab:gates}); the eighth is a wall-clock budget missed only
under transient machine load (nominal projection is within budget). The
anchor result is the dual variable: earlier configurations saturated
$\lambda$ at its cap $\lambda_{\max}=100$ (forcing an always-retrain
degenerate policy); the data-informed target violation
in~\eqref{eq:dual} equilibrates $\lambda$ at a mean of $11.2$. The policy
is non-collapsed (max action share $0.73<0.85$), exploring (entropy $0.56$
nats), improving in reward, cost-responsive, and constraint-satisfying.
This diagnostic does not by itself prove a Pareto win; it proves the
substrate on which the confirmatory run stands is sound---a check we
believe should be standard.

\begin{table}[t]
\centering
\caption{Pre-registered Stage-1 policy-health diagnostic (3 seeds,
independent policies). 7/8 pass; G8 is a transient-load wall-clock miss.}
\label{tab:gates}
\setlength{\tabcolsep}{4pt}
\begin{tabular}{@{}llc@{}}
\toprule
\textbf{Gate} & \textbf{Criterion} & \textbf{Result}\\
\midrule
G1 no collapse       & max action share $\le 0.85$      & $0.73$ \yes\\
G2 exploration       & entropy $\ge 0.15$ nats          & $0.56$ \yes\\
G3 reward improves   & last vs first third              & $+0.029$ \yes\\
G4 cost is live      & policy engages cost trade-off    & \yes\\
G5 no SLA-gaming     & both penalties shrink            & \yes\\
G6 dual $\lambda\le50$ & mean of last third             & $11.2$ \yes\\
G7 policy indep.     & independent per-seed checkpoints & 3/3 \yes\\
G8 wall $\le 6$h     & total across seeds               & $6.8$h $\triangle$\\
\bottomrule
\end{tabular}
\end{table}

\subsection{Is PPO necessary? A policy comparison}\label{sec:ppo-necessity}
The action space is only $\{\textsc{no-op},\textsc{partial},\textsc{full}\}$,
so a reviewer rightly asks whether reinforcement learning is warranted. We
evaluated five policies on the identical env / reward / cost / SLA over the
three contested scenarios ($n{=}9$ cells; Table~\ref{tab:ppo}): no-op,
always-partial, always-full, a hand greedy rule (using the same
attribution-concentration signal the RSO sees), and the learned PPO.
\textbf{The result is a null for the strong claim: PPO is not demonstrably
superior.} On the composite reward, always-partial is best and PPO is
mid-pack; on F1, the greedy rule and always-full lead; and PPO's reward edge
over the greedy rule is not significant (paired Wilcoxon $p=0.29$). PPO does
learn a cost-aware middle policy---less compute than always/greedy-full at
F1 near always-partial---but ``cost-aware middle'' is not ``necessary.'' We
therefore do not claim PPO superiority for this discrete-3 setting; the value
of the RL formulation is expected to appear with a larger or continuous
action space (fraction-of-layers scope), which is future work. This is
consistent with the bang-bang collapse observed under short-horizon training.

\begin{table}[t]
\centering
\caption{Policy comparison on the contested scenarios ($n{=}9$; same
env/reward/cost/SLA). PPO is competitive but not superior; always-partial
wins on reward, greedy/always-full on F1. PPO vs greedy reward: $p=0.29$.}
\label{tab:ppo}
\setlength{\tabcolsep}{4pt}
\begin{tabular}{@{}lcccc@{}}
\toprule
\textbf{Policy} & \textbf{mean F1} & \textbf{GPU-hr} & \textbf{SLA viol.} & \textbf{reward}\\
\midrule
no-op            & $0.564$ & $0$        & $2.33$ & $-0.137$\\
always-partial   & $0.590$ & $1.7\mathrm{e}{-}8$ & $2.33$ & $\mathbf{-0.118}$\\
always-full      & $0.611$ & $3.7\mathrm{e}{-}6$ & $2.00$ & $-1.058$\\
greedy rule      & $\mathbf{0.611}$ & $3.2\mathrm{e}{-}6$ & $2.00$ & $-0.749$\\
\textbf{PPO (ours)} & $0.588$ & $1.2\mathrm{e}{-}6$ & $2.22$ & $-0.438$\\
\bottomrule
\end{tabular}
\end{table}

\subsection{Failure modes}\label{sec:failures}
We surface failures explicitly. (1)~\emph{PSI saturation}: on easy
scenarios PSI is already perfect and CADENCE cannot improve on it.
(2)~\emph{Unrecoverable concept flip}: when the labeling rule inverts, no
covariate-shift retrain recovers F1; the surrogate over-predicts and the
guard must halt retraining. (3)~\emph{Weak encoders mask drift}: on
hash-encoded Airlines and a year-split of Yelp, the trigger never fires
and ``do nothing'' trivially wins---a property of the encoder, not the
loop. (4)~\emph{Single-task forgetting}: EWC-partial has no advantage when
there is no distinct task to preserve. Each of these is a measured
finding, not a conjecture.

%=====================================================================
\section{Limitations and Pre-registered Confirmatory Agenda}\label{sec:limits}
%=====================================================================
We list the limitations a careful reviewer will (rightly) raise, and for
each we state whether it bounds a \emph{claim} or is \emph{planned,
pre-registered} work. We prefer to name these than to have them found.

\begin{itemize}
\item \textbf{Causal validity is shown on synthetic drift only.} Our
ground-truth root causes come from a synthetic intervention mechanism;
real-drift experiments are graded on outcome, not on recovered cause.
We do \emph{not} claim causal identification in the wild. The planned
cross-distribution test trains attribution on intervention mechanisms
$\{A,B,C\}$ and evaluates on unseen $\{D,E,F\}$.
\item \textbf{H2 is PRACTICAL, not STRONG.} The Elec2 result is $n{=}3$
with a deterministic rule policy. The pre-registered confirmatory run is
$10$ seeds $\times$ $30$k timesteps $\times$ $8$ scenarios with the learned
policy, graded by paired non-inferiority; its protocol and success
criteria are fixed in advance.
\item \textbf{PPO necessity is measured---and not established.} We ran the
comparison (Sec.~\ref{sec:ppo-necessity}): against no-op, always-partial,
always-full and a greedy rule, PPO is competitive but not superior
($p=0.29$ vs greedy). We do not claim PPO is necessary at this scale; a
contextual-bandit / supervised-oracle arm and a larger action space are the
remaining tests.
\item \textbf{Causal targeting does not drive the H3 forgetting win (measured null).}
We ran the isolation ablation (Sec.~\ref{sec:h3}): under identical EWC,
CDAG-guided targeting is statistically tied with random/gradient targeting on
forgetting. We therefore claim only EWC-partial $<$ full, not that causal
targeting reduces forgetting. A harder task pair and the learned scorer may
yet reveal a targeting effect.
\item \textbf{Surrogate sample-efficiency: measured; OOD split still open.}
The error-vs-budget curve is now measured (Sec.~\ref{sec:surrogate}):
$\sim\!50$--$100$ interventions reach $R^2\!\approx\!0.9$. The
seen-vs-unseen intervention-type generalization remains to be run.
\item \textbf{End-to-end overhead is accounted, and bounds the claim
(measured).} Sec.~\ref{sec:overhead} shows CADENCE's overhead is dominated by
CDAG construction ($374$\,ms/window) and that, on a trivially small model,
the framework is \emph{net-negative} on compute; net savings require a
full-retrain cost above $\approx1.3$\,s. The saving is thus a
retraining-compute saving that becomes a net system saving only in the
expensive-retrain regime. Confirming the crossover with a direct large-model
retrain measurement is planned.
\item \textbf{Carbon is model-derived, not measured.} Cost/carbon use a
calibrated hardware model, not metered kWh. We therefore keep the headline
claim on \emph{compute} (GPU-hours) and treat carbon as a proportional
estimate~\cite{lacoste2019quantifying,patterson2021carbon}.
\item \textbf{Graph construction scales poorly (measured).} We timed
PC~+~NOTEARS as node count grows: $0.30$\,s at $30$ nodes, $5.2$\,s at
$100$, $127$\,s at $300$ ($\approx O(n^{2.6})$), exceeding a $10$-min budget
at $500$. CDAG construction is impractical beyond $\sim\!100$ nodes and is
the dominant system cost (consistent with Sec.~\ref{sec:overhead}). We
therefore build the graph only \emph{on a trigger alert}
(``online-triggered,'' not continuous discovery), and sparser/approximate
discovery or node-count capping is required before any ``online at scale''
claim. This is a genuine bound on high-dimensional deployment.
\item \textbf{Graph recovery is weak on short windows.} Our structural
Hamming proxy indicates the recovered CDAG under-connects; the GNN wins
\emph{despite} this via message passing, which is why we claim
attribution/prediction (B/C), not recovery (A).
\end{itemize}

None of these undermine the established claims (H1 contested, H3
multi-task, surrogate calibration, robustness, generality, policy health);
they bound the \emph{H2-STRONG} and \emph{causal-in-the-wild} claims,
which we deliberately do not assert.

%=====================================================================
\section{Conclusion}\label{sec:conclusion}
%=====================================================================
CADENCE reframes model maintenance from ``when to retrain'' to
``minimum-cost repair under a performance constraint,'' and instantiates
that reframing as a single closed loop: a causal drift-attribution graph
over external \emph{and} internal nodes, a counterfactual surrogate that
predicts repair benefit, and a cost-constrained policy formalized as a
CMDP. The technical core is established with pre-registered, reproducible
evidence: learned attribution beats correlation where the problem is
actually hard ($p<0.001$); targeted partial retraining eliminates
catastrophic forgetting on a genuinely multi-task benchmark against a fair
replay baseline ($p<0.001$); the counterfactual surrogate is
well-calibrated ($R^2=0.89$); attribution degrades gracefully under
missing telemetry; and---uncommonly for RL-for-systems---we supply a
falsifiable diagnostic that the learned policy is healthy rather than
silently collapsed.

Equally, we are explicit about the frontier, including where our own method
falls short. Two gaps we closed with measurements returned honest
boundaries rather than wins: full system-overhead accounting shows the
cost saving is a \emph{retraining}-compute saving that is net-positive only
above a measured break-even retrain cost ($\approx1.3$\,s), and the
causal-targeting isolation returns a null---forgetting reduction is owed to
EWC, not to causal layer selection, on our benchmark. A third measured boundary: PPO is competitive but \emph{not} superior to a
greedy rule on the discrete-three action space ($p=0.29$), so we do not
claim the RL policy is necessary at this scale. The cost-saving result
itself is a reproducible Pareto \emph{practical} win ($69\%$ retraining
compute for $1.7$ F1 points on real drift); its strict-dominance form and
out-of-distribution causal generalization remain the subject of a
\emph{pre-registered} confirmatory agenda rather than of over-claimed
results. We contend that this separation---strong claims where the
evidence is strong, published protocols where it is not yet---is what makes
the system both scientifically defensible and deployable. For the
failure-prone core---a cost-aware policy that does not collapse---the
proposed constrained-RL training procedure is shown to avoid the principal
policy-degeneracy modes observed during development (bang-bang collapse and
dual saturation); what remains is to establish that the learned policy is
not merely healthy but \emph{necessary and superior} to simpler baselines,
and to confirm the cost saving at scale.

%=====================================================================
\section*{Reproducibility}
All quantitative claims are generated from a measured results ledger; every
experiment lists its exact command, seeds, and configuration hash. The
Stage-1/Stage-2 pre-registration, gate definitions, and the confirmatory
run protocol are fixed in advance and versioned in the repository. Figures
regenerate from experiment artifacts via a single script.

\bibliographystyle{IEEEtran}
\bibliography{references}

\end{document}
